\documentclass[11pt]{article}

\usepackage[margin=1in]{geometry}
\usepackage{amsmath,amssymb,amsfonts}
\usepackage{booktabs}
\usepackage{hyperref}
\usepackage{microtype}
\usepackage{xcolor}

\title{Budgeted Long-Horizon Risk Control via Cost-to-Go Guided Policy Improvement}
\author{Julian Quick}
\date{\today}

\newcommand{\E}{\mathbb{E}}
\newcommand{\R}{\mathbb{R}}
\newcommand{\Qr}{Q_{\mathrm r}}
\newcommand{\Qc}{Q_{\mathrm c}}
\newcommand{\Ct}{C_t}
\newcommand{\Bt}{B_t}

\begin{document}
\maketitle

\begin{abstract}
Many real control constraints are neither instantaneous nor fixed at training time.
Safety contacts in navigation and fatigue loads in wind farms are accumulated over a deployment horizon, and the effect of a risky action may appear only after the state evolves.
For wind farms this delay is physical: wakes convect downstream, so a yaw decision can change future loading on other turbines after a wind-speed-dependent delay.
This draft proposes a stronger alternative to post-hoc instantaneous penalty scheduling: learn long-horizon reward and cost-to-go critics, then perform test-time constrained policy improvement under an explicit remaining budget.
The central object is $\Qc(s,h,a)$, the expected future budget consumption after action $a$ from physical state $s$ and history $h$.
At deployment, a frozen policy proposes actions, while a planner or energy optimizer selects actions that maximize $\Qr$ subject to remaining budget feasibility under $\Qc$.
The method targets unsafe-budget Safety Gym and long-horizon wind-farm load constraints with the same mechanism.
\end{abstract}

\section{Motivation}

The previous post-hoc budget-scheduling paradigm is too weak for the real wind-farm problem.
It modulates an inference-time penalty as the remaining budget changes, but it does not by itself reason about delayed state-mediated cost.
This distinction is decisive.
In Safety Gym, action changes position and velocity; future position determines hazard exposure; hazard exposure consumes safety budget.
In wind-farm control, yaw changes wake structure; wakes convect; downstream turbines experience delayed changes in loading and power.
In both domains, the relevant question is not ``is this action locally unsafe?'' but:
\begin{quote}
If I take this action now, how much scarce future risk or load budget will I consume, and is the reward worth it?
\end{quote}

The correct control object is therefore a cost-to-go model, not a one-step cost surrogate.

\section{Problem Formulation}

We consider an episodic control problem with horizon $T$.
The agent observes a physical state $s_t$, a history summary $h_t$, remaining time $T-t$, and cumulative budget consumption $\Ct$.
It selects a continuous action $a_t \in \mathcal A$.
The environment returns reward $r_t$ and cost $c_t \ge 0$.
The deployment objective is
\begin{align}
  \max_\pi \quad & \E_\pi \left[\sum_{t=0}^{T-1} r_t \right] \\
  \text{s.t.} \quad & \sum_{t=0}^{T-1} c_t \le B
  \quad \text{per episode, or per controlled asset.}
\end{align}

The important setting is state-mediated cost:
\begin{equation}
  c_{t+k} = f(s_{t+k}, h_{t+k}, a_{t+k}),
  \qquad
  s_{t+k} \sim p(\cdot \mid s_t, h_t, a_t, \ldots).
\end{equation}
Thus the consequence of $a_t$ may not be visible in $c_t$.

\paragraph{Wind-farm instance.}
The action is turbine yaw.
The cost is fatigue or load budget consumption.
The history $h_t$ must summarize yaw history, wind speed, wind direction, turbulence, and pending wake arrivals.
The delayed effect is approximately governed by convective travel time
\begin{equation}
  \tau_{ij}(t) \approx \frac{d_{ij}^{\parallel}(t)}{u_t},
\end{equation}
where $d_{ij}^{\parallel}$ is downstream spacing in the wind-aligned frame and $u_t$ is wind speed.

\paragraph{Safety Gym instance.}
The action changes robot motion.
The cost is unsafe contact or hazard exposure.
The history may be short because the simulator state is nearly Markov, but the action-cost link is still mediated by dynamics.

\section{Method}

\subsection{Learn long-horizon critics}

Given trajectories from an unconstrained or weakly constrained policy, fit
\begin{align}
  \Qr(s_t,h_t,a_t) &\approx \E\left[\sum_{k=0}^{T-t-1} \gamma_r^k r_{t+k}\mid s_t,h_t,a_t\right],\\
  \Qc(s_t,h_t,a_t) &\approx \E\left[\sum_{k=0}^{T-t-1} \gamma_c^k c_{t+k}\mid s_t,h_t,a_t\right].
\end{align}

For safety, $\Qc$ should be conservative.
Practical variants include twin critics with max aggregation, ensembles with mean-plus-variance pessimism, and distributional critics estimating a high quantile of future cost.

\subsection{Test-time constrained policy improvement}

At deployment, a frozen actor supplies candidate actions.
The controller improves those candidates using the learned critics.
The hard form is
\begin{equation}
  a_t^\star =
  \arg\max_{a \in \mathcal A}
  \Qr(s_t,h_t,a)
  \quad
  \text{s.t.}
  \quad
  \Ct + \Qc(s_t,h_t,a) \le B.
\end{equation}

The soft form is
\begin{equation}
  a_t^\star =
  \arg\max_{a \in \mathcal A}
  \Qr(s_t,h_t,a)
  - \lambda_t \Qc(s_t,h_t,a)
  - \beta \sigma_{\Qc}(s_t,h_t,a),
\end{equation}
where $\lambda_t$ is a budget shadow price and $\sigma_{\Qc}$ is epistemic uncertainty.
The old urgency schedule can initialize $\lambda_t$, but it is no longer the main contribution.

\subsection{Candidate selection and energy optimization}

There are two implementation paths.
For Gaussian SAC actors, sample $K$ candidate actions around the frozen actor and select the best feasible action.
For energy-based actors, add $\lambda_t \Qc$ to the action energy and optimize actions by gradient descent.
Both instantiate the same constrained policy-improvement step.

\section{Why this is stronger than instantaneous penalties}

Instantaneous penalties can only discourage actions that look unsafe now.
They fail when the risk is delayed.
Cost-to-go guidance can rank an action as expensive because it creates future hazard exposure or future wake-induced load.
This is exactly the structural issue shared by Safety Gym and wind-farm control.

\begin{table}[h]
\centering
\begin{tabular}{lll}
\toprule
Domain & Delayed mechanism & Required critic input \\
\midrule
Safety Gym & action $\rightarrow$ motion $\rightarrow$ hazard exposure & simulator state, velocity, action \\
Wind farm & yaw $\rightarrow$ wake convection $\rightarrow$ downstream load & wind, geometry, yaw history, pending wakes \\
\bottomrule
\end{tabular}
\caption{Both target domains require future-cost reasoning rather than local action penalties.}
\end{table}

\section{Research Bet and Kill Tests}

This project should be pursued only in a narrow form.
The contribution is not ``learn a safety critic for constrained RL,'' which is already close to existing work on constrained policy optimization, budgeted RL, Lyapunov methods, and distributional safety critics.
The contribution must be sharper:
\begin{quote}
Test-time constrained improvement over a frozen policy can solve delayed budget spending better than instantaneous penalties, static schedules, or retraining baselines, especially when the cost is state-mediated by dynamics such as hazard exposure or wake convection.
\end{quote}

\paragraph{Nugget.}
Delayed constraints should be controlled as scarce future-budget allocation problems, not as instantaneous penalty problems.

\paragraph{Main technical risk.}
The learned cost-to-go critic may rank actions correctly under the data-collection policy but become miscalibrated under the improved closed-loop controller.
This is the core distribution-shift risk.
The project should therefore test ranking and calibration before investing in paper-scale baselines.

\paragraph{First question to resolve.}
Can $\Qc$ rank candidate actions by true future cost under the closed-loop improved controller, not merely under the behavior policy used to collect training data?

\subsection{Critical experiment sequence}

\paragraph{Experiment 1: oracle delayed-cost sanity check.}
In the wake-convection proxy, compare oracle future-cost candidate selection against one-step penalties and the closed-form urgency schedule.
The oracle estimates the future load consequence of a candidate by rolling out the proxy dynamics.
This isolates whether future-cost information is actually useful.
\emph{Kill criterion:} if oracle future-cost selection does not improve the reward--budget Pareto frontier, the core premise is weak.

\paragraph{Experiment 2: learned critic ranking and calibration.}
Train $\Qc$ on frozen-policy trajectories, then evaluate whether it ranks candidate actions by true rollout cost.
Report rank correlation, calibration curves, and closed-loop budget satisfaction.
\emph{Kill or park criterion:} if $\Qc$ cannot rank actions in the cheap wake proxy with adequate data, the method is not ready for Safety Gym or wind farms.

\paragraph{Experiment 3: unsafe-budget Safety Gym.}
Compare frozen actor, one-step penalty, urgency schedule, SAC-Lagrangian or CPO-style retraining, and $\Qc$ test-time improvement across budgets.
The target is not zero cost; the target is high reward with high budget satisfaction and meaningful budget utilization.
\emph{Kill criterion:} if $\Qc$ only wins by becoming conservative and underusing the unsafe budget, the Safety Gym claim should be dropped or reframed.

\section{Experimental Plan}

\subsection{Testbed 1: Wake-convection proxy}

Before the real loading surrogate is available, use a deliberately cheap proxy with:
\begin{itemize}
  \item wind-speed-dependent wake delays between turbine pairs,
  \item action history buffers representing pending wake arrivals,
  \item immediate yaw-rate/self-load costs,
  \item delayed downstream load costs caused by upstream yaw,
  \item time-varying reward opportunities so spending budget can be worthwhile.
\end{itemize}

This proxy is not intended to be physically publishable.
Its job is to test whether the algorithm solves the delayed-budget structure.

\subsection{Testbed 2: Safety Gym unsafe budget}

Safety Gym should be treated as a hard surrogate, not a weak side experiment.
The target metric is not zero cost.
The target is strategic use of a limited unsafe budget:
\begin{equation}
  \text{high reward}, \qquad \sum_t c_t \le B, \qquad \text{high budget utilization when risk is valuable.}
\end{equation}

\subsection{Baselines}

\begin{enumerate}
  \item Unconstrained frozen actor.
  \item Hard safety filter or action clipping.
  \item One-step cost penalty.
  \item Closed-form urgency schedule without $\Qc$.
  \item SAC-Lagrangian retraining for each budget.
  \item Model-predictive control with learned dynamics and learned cost.
\end{enumerate}

\subsection{Metrics}

Report per-episode budget satisfaction, mean and lower-tail reward, budget utilization, violation severity, critic calibration, critic ranking accuracy, and generalization across budgets.
For wind farms, additionally report power retention, per-turbine load budget satisfaction, and zero-shot layout transfer.

\section{Expected Failure Modes}

\paragraph{Cost critic miscalibration.}
The controller may satisfy the learned budget but violate the true budget.
This motivates conservative aggregation, calibration curves, and online correction.

\paragraph{Distribution shift.}
The improved controller may choose actions outside the offline data support.
This motivates ensemble pessimism, behavior-closeness penalties, and candidate generation from the frozen actor rather than arbitrary random search.

\paragraph{Overconservatism.}
A pessimistic $\Qc$ can avoid all risk and underuse budget.
Budget utilization should be a primary metric, not an afterthought.

\section{Positioning}

The contribution is not a new penalty schedule.
The contribution is a deployment-time constrained policy-improvement layer for state-mediated, long-horizon budget constraints.
Safety Gym and wind-farm loads are two instances of the same problem.
The closed-form schedule remains useful only as a budget shadow-price initializer.

\bibliographystyle{plain}
\bibliography{references}

\end{document}
